% =====================================================================
% Lab 1 - Initial Security Assessment report body.
% NOTE ON FIGURES: each figure below shows a placeholder box so the file
% compiles as-is. To insert a screenshot, drop the image into figures/
% and replace the \fbox{...} line with the commented \includegraphics line
% above it. Keep the version numbers and packet counts in the text
% consistent with whatever screenshots you finally use.
% =====================================================================


% \begin{center}
% \begin{tabular}{ll}
% \textbf{Course:} & NETAUTH -- Network Security: Authentication and Authorization \\
% \textbf{Date:} & 23 September 2026 \\
% \textbf{Scope:} & Isolated laboratory network 192.168.1.0/24 (lab systems only) \\
% \end{tabular}
% \end{center}

% \vspace{0.3cm}

% ---------------------------------------------------------------------
\section{Introduction and Objective}

The goal of this laboratory was to carry out an initial security assessment of
a deliberately insecure internal web service. The assessment was performed from
a Kali Linux virtual machine against an Ubuntu Server virtual machine on the
same isolated network. The work covered network connectivity, service
discovery, exploration of the web application and its access control,
administrator credential testing, and analysis of the login traffic.

All activity was limited to the laboratory systems. The findings and
conclusions below are based only on the evidence collected during the exercise.

% ---------------------------------------------------------------------
\section{Network Topology and Configuration}

The laboratory consisted of two virtual machines on the same \texttt{/24}
network, as shown in Table~\ref{tab:topology}.

\begin{table}[H]
\centering
\caption{Laboratory network configuration.}
\label{tab:topology}
\begin{tabular}{@{}llll@{}}
\toprule
\textbf{System} & \textbf{Role} & \textbf{IP address} & \textbf{Network} \\
\midrule
Kali Linux    & Assessment client   & 192.168.1.5  & 192.168.1.0/24 \\
Ubuntu Server & Target web server   & 192.168.1.35 & 192.168.1.0/24 \\
\bottomrule
\end{tabular}
\end{table}

On the Kali machine, interface \texttt{eth0} was assigned the address
192.168.1.5/24 with a directly connected route to 192.168.1.0/24, which placed
both machines on the same isolated segment.

\begin{figure}[H]
\centering
\includegraphics[width=0.6\textwidth]{figures/fig1_kali_ip.png}
\caption{Kali Linux network configuration and route for 192.168.1.0/24.}
\label{fig:kali-ip}
\end{figure}

% ---------------------------------------------------------------------
\section{Network Connectivity}

Connectivity from Kali to the Ubuntu Server was checked with ICMP echo
requests. The target 192.168.1.35 answered every request with 0\% packet loss,
which confirmed that the two machines could communicate before any further
testing, as shown in Figure~\ref{fig:ping}.

\begin{figure}[H]
\centering
\includegraphics[width=0.85\textwidth]{figures/fig2_ping.png}
\caption{Successful ping test from Kali Linux to 192.168.1.35.}
\label{fig:ping}
\end{figure}

% ---------------------------------------------------------------------
\section{Service Discovery and Attack Surface}

The \texttt{nmap} tool was used to identify the TCP services reachable on the Ubuntu Server. A
version detection scan found two open ports, summarised in
Table~\ref{tab:nmap}.

\begin{table}[H]
\centering
\caption{Open TCP ports found on 192.168.1.35.}
\label{tab:nmap}
\begin{tabular}{@{}llp{8.2cm}@{}}
\toprule
\textbf{Port} & \textbf{Service} & \textbf{Observed version / interpretation} \\
\midrule
22/tcp   & SSH  & OpenSSH 10.2p1 Ubuntu 2ubuntu3.6; remote administration service. \\
5000/tcp & HTTP & Werkzeug httpd 3.1.8 (Python 3.14.4); hosts the laboratory web application. \\
\bottomrule
\end{tabular}
\end{table}

Both services are reachable over the network and therefore add to the attack
surface, as shown in Figure~\ref{fig:nmap}. The scan alone does not prove that SSH is vulnerable; it only shows
that SSH is an additional remotely accessible service that must be securely
configured. Port 5000 exposes the web application examined in the next
sections. 
% It is also worth noting that the web application is served by
% Werkzeug, the development server built into Flask, which is not intended to be
% exposed as a production service and whose version is disclosed in the response.

\begin{figure}[H]
\centering
\includegraphics[width=0.9\textwidth]{figures/fig3_nmap.png}
\caption{Nmap service-version scan against 192.168.1.35.}
\label{fig:nmap}
\end{figure}

% ---------------------------------------------------------------------
\section{Web Application Exploration}

\subsection{Public resources}

Before authentication, the root page at \url{http://192.168.1.35:5000/} reported
that the user was not authenticated and offered links to \emph{Public
information} and \emph{Login}. The landing page and the public information page
were both reachable without logging in, which is the expected behaviour for
public content.

\begin{figure}[H]
\centering
\includegraphics[width=0.85\textwidth]{figures/fig4_landing.png}
\caption{Web application landing page before authentication.}
\label{fig:landing}
\end{figure}

\subsection{Authenticated employee access}

The legitimate account \texttt{alice / alice123} was used to log in. After
authentication, Alice was identified with the \emph{employee} role and could
open \texttt{/internal}, as shown in Figure~\ref{fig:internal}. That page disclosed the internal network 192.168.1.0/24
and stated that the information should only be available to authenticated users.
For an authenticated employee, this access is consistent with the stated policy.

\begin{figure}[H]
\centering
\includegraphics[width=0.85\textwidth]{figures/fig5_internal.png}
\caption{Alice authenticated as an employee and accessing \texttt{/internal}.}
\label{fig:internal}
\end{figure}

\subsection{Unauthorised access to the administration page}

While still logged in as Alice, the \emph{Administration} link opened
\texttt{/admin}, as shown in Figure~\ref{fig:admin-employee}. The page showed Alice's role as \emph{employee} and, at the
same time, stated that it should only be accessible to administrators. This is
direct evidence that the application does not enforce role-based authorisation
on the administrative resource: being authenticated is enough to reach it,
regardless of the user's role.

\begin{figure}[H]
\centering
\includegraphics[width=0.85\textwidth]{figures/fig6_admin_employee.png}
\caption{Employee user Alice accessing the Administration panel.}
\label{fig:admin-employee}
\end{figure}

\subsection{Administrator credential test}

The administrator account was then accessed using the guessed credentials
\texttt{admin / admin123}. The credentials followed the same predictable pattern
as the other accounts and were accepted, giving a session identified as
\emph{admin}. The success of a simple, predictable password shows that
administrator authentication is weak.

\begin{figure}[H]
\centering
\includegraphics[width=0.85\textwidth]{figures/fig7_admin_login.png}
\caption{Successful administrator session using the guessed credentials.}
\label{fig:admin-login}
\end{figure}

% ---------------------------------------------------------------------
\section{Traffic Analysis}

Wireshark was started on the Kali \texttt{eth0} interface, and Alice logged in
again, as shown in Figure~\ref{fig:http-flow}. Filtering the capture for the web traffic showed the login sequence: a
\texttt{POST /login} request, an HTTP 302 redirect, and a following request for
\texttt{/internal}.

\begin{figure}[H]
\centering
\includegraphics[width=0.9\textwidth]{figures/fig8_http_flow.png}
\caption{HTTP traffic showing \texttt{POST /login}, the redirect, and access to \texttt{/internal}.}
\label{fig:http-flow}
\end{figure}

Inspecting the \texttt{POST /login} packet showed that the request used
\texttt{application/x-www-form-urlencoded} over plain HTTP. The form fields were
readable directly in Wireshark: \texttt{username = alice} and
\texttt{password = alice123}, as shown in Figure~\ref{fig:creds}. The login credentials were therefore sent over the
network without any transport encryption.

\begin{figure}[H]
\centering
\includegraphics[width=0.9\textwidth]{figures/fig9_wireshark_creds.png}
\caption{Wireshark packet details exposing Alice's username and password in the HTTP request.}
\label{fig:creds}
\end{figure}

% ---------------------------------------------------------------------
\section{Security Assessment}

Each finding below follows the required \emph{Observation $\rightarrow$ Security
consequence $\rightarrow$ Countermeasure} structure and states the security
properties affected. Only issues supported by the collected evidence are
included.

\subsection{Broken access control on \texttt{/admin}}
\textbf{Observation.} Alice is authenticated with the \emph{employee} role but
can open the \texttt{/admin} page, which itself states that it should be
restricted to administrators.\\
\textbf{Security consequence.} A non-administrator can reach administrative
functionality and information. If the listed administrative functions were
active, this could expose sensitive data or allow unauthorised configuration
changes.\\
\textbf{Affected properties.} Authorization; Confidentiality; potentially
Integrity.\\
\textbf{Countermeasure.} Enforce server side role checks on every
administrative endpoint. Authentication alone must not grant access; the server
must confirm that the user holds the \emph{admin} role before returning the page
or performing any privileged action.

\subsection{Weak and predictable administrator credentials}
\textbf{Observation.} The administrator account was accessed using the guessed
password \texttt{admin / admin123}, which follows the same obvious pattern as
the other accounts.\\
\textbf{Security consequence.} An attacker who guesses the password gains
administrator privileges and access to privileged resources, which undermines
the authentication mechanism.\\
\textbf{Affected properties.} Authentication; Authorization; Confidentiality;
potentially Integrity.\\
\textbf{Countermeasure.} Require strong, unique administrator passwords; forbid
predictable or default-style passwords; add rate limiting or account lockout on
repeated failures; and use multi-factor authentication for privileged accounts
where possible.

\subsection{Credentials transmitted in clear text over HTTP}
\textbf{Observation.} Wireshark captured the \texttt{POST /login} request and
showed the fields \texttt{username = alice} and \texttt{password = alice123} in
readable form. The application was accessed over HTTP, not HTTPS.\\
\textbf{Security consequence.} Anyone able to observe the traffic can recover
user credentials and reuse them to impersonate the user. Other HTTP content is
also open to passive observation or modification in transit.\\
\textbf{Affected properties.} Confidentiality; Authentication; potentially
Integrity of data in transit.\\
\textbf{Countermeasure.} Serve the application only over HTTPS with TLS,
redirect HTTP to HTTPS, and apply secure session-cookie settings once TLS is in
place.

\subsection{Exposed network services increase the attack surface}
\textbf{Observation.} Nmap found SSH on TCP/22 and the web service on TCP/5000
reachable from the Kali host, and both responses disclosed their software
versions.\\
\textbf{Security consequence.} Each reachable service is an additional entry
point. The scan alone does not prove an exploitable flaw, but unnecessary or
weakly configured services increase the opportunities available to an attacker.\\
\textbf{Affected properties.} No direct violation was shown by the scan itself;
a compromise of either service could affect Confidentiality, Integrity, or
Availability.\\
\textbf{Countermeasure.} Expose only the services that are required, restrict
access with host or network firewall rules, keep software updated, use strong
authentication for SSH, and avoid exposing development web servers beyond the
intended trusted network.

\subsection{Summary of affected security properties}

Table~\ref{tab:matrix} maps each finding to the security properties it affects.

\begin{table}[H]
\centering
\caption{Findings and affected security properties.}
\label{tab:matrix}
\small
\begin{tabularx}{\textwidth}{@{}Xccccc@{}}
\toprule
\textbf{Finding} & \textbf{Conf.} & \textbf{Integ.} & \textbf{Avail.} & \textbf{Authn.} & \textbf{Authz.} \\
\midrule
Broken access control          & Yes       & Potential & -- & --        & Yes \\
Weak admin credentials         & Yes       & Potential & -- & Yes       & Yes \\
Plaintext HTTP credentials     & Yes       & Potential & -- & Yes       & --  \\
Reachable SSH + HTTP services  & Potential & Potential & Potential & Potential & Potential \\
\bottomrule
\end{tabularx}
\end{table}

% ---------------------------------------------------------------------
\section{Conclusion}

The initial assessment shows that the service is reachable and working but has
several important weaknesses. The most significant confirmed issues are broken
authorisation on the administration page, a predictable administrator password,
and login credentials sent in clear text over HTTP. Together, these problems let
a low-privilege user reach administrator-only content, make compromise of the
administrator account easy, and expose user credentials to anyone capturing the
network traffic.

The first priorities for remediation should be to enforce server-side
authorisation on administrative resources, replace the weak administrator
password with strong authentication controls, and serve all traffic over
HTTPS/TLS. The exposed SSH and web services should also be hardened and limited
to the systems that need them. Availability was not affected during the
laboratory activities.

% ---------------------------------------------------------------------
\appendix
\section{Evidence Checklist}

Table~\ref{tab:checklist} maps each required deliverable to its location in this
report.

\begin{table}[H]
\centering
\caption{Required deliverables and where they appear.}
\label{tab:checklist}
\begin{tabularx}{\textwidth}{@{}Xl@{}}
\toprule
\textbf{Required deliverable} & \textbf{Location} \\
\midrule
Network topology and configuration            & Section 2, Fig.~\ref{fig:kali-ip} \\
Connectivity evidence                         & Section 3, Fig.~\ref{fig:ping} \\
Reconnaissance results and interpretation     & Section 4, Fig.~\ref{fig:nmap} \\
Resources discovered in the web application    & Section 5, Figs.~\ref{fig:landing}--\ref{fig:admin-login} \\
Wireshark evidence of the authentication exchange & Section 6, Figs.~\ref{fig:http-flow}--\ref{fig:creds} \\
Identified vulnerabilities and affected properties & Section 7 \\
Proposed countermeasures                      & Section 7 \\
Security-posture conclusion                    & Section 8 \\
\bottomrule
\end{tabularx}
\end{table}
